\documentclass[11pt]{article}
\usepackage[a4paper,margin=1in]{geometry}
\usepackage{amsmath,amssymb}
\usepackage{booktabs}
\usepackage[numbers,sort&compress]{natbib}
\usepackage{hyperref}
\hypersetup{colorlinks=true,linkcolor=blue,citecolor=blue,urlcolor=blue}
\newcommand{\doi}[1]{DOI: \href{https://doi.org/#1}{#1}}

\title{Uniform Rotation is the Field-Cost Optimum over a Domain-Wall Sweep:\\
A First Step Beyond the Macrospin in Optimal Magnetization Switching}
\author{Felipe Santib\'a\~nez-Leal\\
\small CAOS open-research programme, Santiago, Chile}
\date{2026}

\begin{document}
\maketitle

\begin{abstract}
Optimal control of magnetization switching has been solved at the macrospin level, where the magnetic
element is treated as a single rotating moment. The authors of the method state, in print, that this
approximation breaks down with size and that the reversal may instead proceed by nonuniform rotation,
domain-wall nucleation and propagation, or spin waves, and that it remains open under what conditions
these become optimal in terms of energy efficiency. We take a first step past the macrospin. For a
ferromagnetic spin chain with nearest-neighbour exchange, we compute the switching cost of two
reversal modes, a uniform rotation in which every site follows the macrospin optimal control path
together, and a domain-wall sweep in which a reversed domain nucleates at one end and its wall
propagates to the other. The lattice engine reproduces the analytic macrospin cost in the one-site
limit. We find that for the switching-cost metric, the Joule heating of the source, uniform rotation
is cheaper than the domain-wall sweep across every chain length from two to sixty-four sites and every
exchange strength tested, because the wall forces fast local flips and bends the bonds across itself
while uniform rotation moves every site slowly and never pays exchange. This is worth stating plainly
because it is the opposite of a common intuition: domain-wall motion dominates real, thermally-driven
switching, but that is a thermal-barrier advantage, not a field-cost advantage. All results are
reproducible from committed artifacts and a live instance.
\end{abstract}

\section{Introduction}

Optimal control theory computes the applied-field pulse that reverses a magnetic moment for the least
dissipated energy~\citep{kwiatkowski2021}, and applied to two-dimensional van der Waals magnets it
reaches picosecond, energy-efficient switching~\citep{badarneh2026}. The formalism has been developed
almost entirely at the macrospin level, a single moment on the unit sphere. Its own authors note the
limitation explicitly: ``A macrospin approximation is used in the present study, but this is expected
to break down with increasing system size. Even if the initial and the final states are collinear, the
transition between them may involve nonuniform rotation of magnetization such as nucleation and
propagation of domain walls or excitation of spin waves. It remains to be seen under what conditions
these and possibly other, yet unknown switching mechanisms become optimal in terms of energy
efficiency''~\citep{badarneh2023}.

The only step past the macrospin taken anywhere is a single one-dimensional study, which found that
above a critical length a standing spin wave appears in the optimal reversal~\citep{badarneh2020}. Here
we compare, for a spin chain with exchange, the two reversal modes that matter most for a device, a
uniform rotation and a domain-wall sweep, on the switching-cost metric, and we identify which is
cheaper and why.

\section{Model}

A chain of $N$ unit moments has energy
\begin{equation}
E = -K \sum_i s_{i,z}^2 \;-\; J \sum_{\langle ij\rangle} \vec s_i \cdot \vec s_j,
\end{equation}
with $K$ the easy-axis anisotropy per site and $J$ the nearest-neighbour exchange (positive is
ferromagnetic), open boundaries. The internal field at a site adds an exchange term
$(J/\mu)(\vec s_{i-1} + \vec s_{i+1})$ to the single-site anisotropy field. In the one-site,
no-exchange limit the chain is a macrospin, and our lattice cost reproduces the closed-form macrospin
cost to within $2\times10^{-3}$, which is the validation gate for the method.

The switching cost of any chain trajectory is computed as everywhere in this framework: the equation of
motion is inverted at every site to give the applied field that site needs, with the chain internal
field (including exchange) subtracted, and the summed squared field is integrated over time,
\begin{equation}
\Phi = \int_0^T \sum_i |\vec b_i(t)|^2 \, dt,
\end{equation}
in units of T$^2$~s. This is not an energy; converting to joules requires a circuit model, which we do
not assume.

\subsection{Two reversal modes}

In the \emph{uniform rotation}, every site follows the same macrospin optimal control path together.
Neighbouring spins stay parallel throughout, so exchange is never paid, and the cost scales linearly in
$N$. In the \emph{domain-wall sweep}, a reversed domain nucleates at one end and a wall of finite width
propagates to the other, so at any instant the chain is part up, part down. Spins across the wall are
not parallel, so exchange is paid, and the wall must cross the chain in the switching time, which forces
the sites it passes to flip quickly.

\section{Results}

Table~\ref{tab:lattice} reports the two costs for a CrSBr chain at exchange $J/K = 0.5$ and switching
time $T = 10\,\tau_0$, as the chain grows from two to sixty-four sites.

\begin{table}[h]
\centering
\caption{Switching cost of the uniform rotation and the domain-wall sweep for a CrSBr chain,
$J/K = 0.5$, $T = 10\,\tau_0$. Uniform rotation is cheaper at every length; its cost scales linearly in
$N$ while the wall cost scales faster.}
\label{tab:lattice}
\begin{tabular}{cccc}
\toprule
$N$ sites & uniform cost (T$^2$~s) & wall cost (T$^2$~s) & ratio wall/uniform \\
\midrule
2  & $1.9\times10^{-11}$ & $2.7\times10^{-10}$ & 13.8 \\
4  & $3.9\times10^{-11}$ & $3.6\times10^{-10}$ & 9.4 \\
8  & $7.8\times10^{-11}$ & $5.9\times10^{-10}$ & 7.6 \\
16 & $1.6\times10^{-10}$ & $1.3\times10^{-9}$  & 8.4 \\
32 & $3.1\times10^{-10}$ & $3.9\times10^{-9}$  & 12.5 \\
64 & $6.2\times10^{-10}$ & $1.4\times10^{-8}$  & 22.0 \\
\bottomrule
\end{tabular}
\end{table}

Uniform rotation is cheaper at every chain length, by a factor between eight and twenty-two. Its cost
doubles as the chain doubles, the linear scaling expected when exchange is never paid. The domain-wall
cost grows faster than linearly, because the wall must sweep a longer chain in the same time and
therefore forces faster local flips, whose field cost grows with their speed. We repeated the
comparison across exchange strengths from $J/K = 0.2$ to $10$ and found uniform rotation cheaper
throughout.

\section{Discussion}

This result is the opposite of a common intuition, and the resolution is instructive. Domain-wall
motion dominates real magnetic switching, and much device engineering is aimed at nucleating and moving
walls at low current. But that dominance is a \emph{thermal-barrier} advantage: a wall lowers the
energy barrier for reversal, so it is easier to trigger by thermal activation or a small current. The
switching cost minimized by optimal control is a different quantity, the Joule heating of the driving
circuit, and for that quantity the cheap thing to do is to move every moment slowly and uniformly,
never bending a bond. The two metrics reward opposite mechanisms.

The comparison here is between two specific ans\"atze, a uniform rotation and a constant-speed
domain wall, not a free search over all chain trajectories. A full lattice optimal control path, the
image-based direct-minimization method generalized to the chain with a per-site field, could in
principle find a cheaper nonuniform mode, for instance one that lets the wall slow down where it can.
That free lattice solve is the natural next step. What the present comparison establishes is the
validation limit, the macrospin reproduction, and a clear, quantitative answer for the two modes that
matter most to a device: for the field-cost metric, uniform rotation wins.

\section{Conclusion}

We have taken a first step past the macrospin in optimal magnetization switching, comparing the
switching cost of a uniform rotation and a domain-wall sweep for a ferromagnetic spin chain. Uniform
rotation is the field-cost optimum across every chain length and exchange strength tested. The
domain-wall mode that dominates real thermally-driven switching is more expensive in field cost, a
distinction between the thermal-barrier and the Joule-heating metrics that this comparison makes
concrete. Results are reproducible from committed artifacts and a live instance at
\url{https://espira.fasl-work.com}.

\section*{Data and code availability}

The engine \texttt{spinoct} is open source (MIT) at
\url{https://github.com/fsantibanezleal/CAOS_SpinOCT}; the product, the bake and the web application at
\url{https://github.com/fsantibanezleal/CAOS_RES_Espira}. Every number in Table~\ref{tab:lattice} is
reproduced from a committed artifact.

\begin{thebibliography}{9}
\bibitem{kwiatkowski2021} G. J. Kwiatkowski, M. H. A. Badarneh, D. V. Berkov, P. F. Bessarab,
\emph{Optimal Control of Magnetization Reversal in a Monodomain Particle by Means of Applied Magnetic
Field}, Phys. Rev. Lett. \textbf{126}, 177206 (2021). \doi{10.1103/PhysRevLett.126.177206}.
\bibitem{badarneh2026} M. H. Badarneh, P. Cai, E. J. G. Santos,
\emph{Optimal Control Drives Ultrafast and Energy-Efficient Magnetization Switching in Van der Waals
Magnets}, Advanced Materials \textbf{e23059} (2026). \doi{10.1002/adma.202523059}.
\bibitem{badarneh2023} M. H. A. Badarneh, G. J. Kwiatkowski, P. F. Bessarab,
\emph{Reduction of energy cost of magnetization switching in a biaxial nanoparticle by use of internal
dynamics}, Phys. Rev. B \textbf{107}, 214448 (2023). \doi{10.1103/PhysRevB.107.214448}.
\bibitem{badarneh2020} M. H. A. Badarneh, G. J. Kwiatkowski, P. F. Bessarab,
\emph{Mechanisms of energy-efficient magnetization switching in a bistable nanowire}, Nanosystems:
Physics, Chemistry, Mathematics \textbf{11}, 294 (2020). \doi{10.17586/2220-8054-2020-11-3-294-300}.
\end{thebibliography}

\end{document}
